% Part: first-order-logic
% Chapter: syntax-and-semantics
% Section: formation-sequences

\documentclass[../../../include/open-logic-section]{subfiles}

\begin{document}

\olfileid{pl}{syn}{fseq}

\olsection{গঠন-অনুক্রম}

আরোহী সংজ্ঞা দিয়ে !!{formula}s সংজ্ঞায়িত করা এবং তার পরিপূরক কৌশল
হিসেবে আরোহের সাহায্যে !!{formula}s-এর ধর্ম প্রমাণ করা—এটি একটি সুন্দর
ও কার্যকর পদ্ধতি। তবে !!{formula}s নির্মাণকে আরও নিচ থেকে উপরে,
ধাপে ধাপে দেখাও উপযোগী হতে পারে। এখানে \emph{গঠন-অনুক্রম}-এর ধারণা
ব্যবহার করে আমরা সেটিই করব।

\begin{defn}[সূত্রের গঠন-অনুক্রম]
\ollabel{defn:fseq-frm}
$\Lang L_0$ ভাষার সংকেত দিয়ে গঠিত প্রতীকক্রমগুলির একটি সসীম অনুক্রম
$\tuple{!A_0,\dotsc,!A_n}$, $!A$-এর জন্য একটি \emph{গঠন-অনুক্রম},
যদি $!A \ident !A_n$ হয় এবং প্রতিটি $i \leq n$-এর জন্য হয় $!A_i$
একটি পরমাণু সূত্র, নয়তো এমন $j,k < i$ আছে যাতে নিচের কোনো একটি সত্য:
\begin{enumerate}
    \tagitem{prvNot}{$!A_i \ident \lnot !A_j$।}{}%
    \tagitem{prvAnd}{$!A_i \ident (!A_j \land !A_k)$।}{}%
    \tagitem{prvOr}{$!A_i \ident (!A_j \lor !A_k)$।}{}%
    \tagitem{prvIf}{$!A_i \ident (!A_j \lif !A_k)$।}{}%
    \tagitem{prvIff}{$!A_i \ident (!A_j \liff !A_k)$।}{}%
\end{enumerate}
\end{defn}

\begin{ex}
\[
\tuple{
    \Obj p_0,
    \Obj p_1,
    (\Obj p_1 \land \Obj p_0),
    \lnot (\Obj p_1 \land \Obj p_0)
}
\]
অনুক্রমটি
$\lnot (\Obj p_1 \land \Obj p_0)$-এর একটি গঠন-অনুক্রম; নিচেরটিও তাই:
\[
\tuple{
    \Obj p_0,
    \Obj p_1,
    \Obj p_0,
    (\Obj p_1 \land \Obj p_0),
    (\Obj p_0 \lif \Obj p_1),
    \lnot (\Obj p_1 \land \Obj p_0)
}.
\]
%
দ্বিতীয় উদাহরণ থেকে দেখা যায়, গঠন-অনুক্রমে “আবর্জনা” থাকতে পারে:
এমন সূত্র, যা অপ্রয়োজনীয় অথবা নির্মাণে কোনো অবদান রাখে না।
\end{ex}

\begin{prop}
\ollabel{prop:formed}
$\Frm[L_0]$-এর প্রতিটি !!{formula}~$!A$-এর একটি গঠন-অনুক্রম আছে।
\end{prop}

\begin{proof}
ধরি $!A$ পরমাণু। তাহলে~$\tuple{!A}$ অনুক্রমটি~$!A$-এর একটি
গঠন-অনুক্রম।
%
এবার ধরি $!B$ ও~$!C$-এর গঠন-অনুক্রম যথাক্রমে
$\tuple{!B_0,\dotsc,!B_n}$ ও $\tuple{!C_0,\dotsc,!C_m}$।
%
\begin{enumerate}
  \tagitem{prvNot}{যদি $!A \ident \lnot !B$ হয়,
      তবে $\tuple{!B_0,\dotsc,!B_n,\lnot !B_n}$
      হলো~$!A$-এর একটি গঠন-অনুক্রম।}{}
  \tagitem{prvAnd}{যদি $!A \ident (!B \land !C)$ হয়,
      তবে $\tuple{!B_0,\dotsc,!B_n,!C_0,\dotsc,!C_m,(!B_n \land !C_m)}$
      হলো~$!A$-এর একটি গঠন-অনুক্রম।}{}
  \tagitem{prvOr}{যদি $!A \ident (!B \lor !C)$ হয়,
      তবে $\tuple{!B_0,\dotsc,!B_n,!C_0,\dotsc,!C_m,(!B_n \lor !C_m)}$
      হলো~$!A$-এর একটি গঠন-অনুক্রম।}{}
  \tagitem{prvIf}{যদি $!A \ident (!B \lif !C)$ হয়,
      তবে $\tuple{!B_0,\dotsc,!B_n,!C_0,\dotsc,!C_m,(!B_n \lif !C_m)}$
      হলো~$!A$-এর একটি গঠন-অনুক্রম।}{}
  \tagitem{prvIff}{যদি $!A \ident (!B \liff !C)$ হয়,
      তবে $\tuple{!B_0,\dotsc,!B_n,!C_0,\dotsc,!C_m,(!B_n \liff !C_m)}$
      হলো~$!A$-এর একটি গঠন-অনুক্রম।}{}
\end{enumerate}
!!{formula}s-এর উপর আরোহের নীতি অনুসারে, প্রতিটি !!{formula}-এর একটি
গঠন-অনুক্রম আছে।
\end{proof}

বিপরীতটিও প্রমাণ করা যায়। এটি গুরুত্বপূর্ণ, কারণ এতে দেখা যায় যে
সূত্র সংজ্ঞায়িত করার আমাদের দুটি উপায় সমতুল্য: দুটিই একই ফল দেয়।
এর আরও একটি অর্থ হলো, গঠন-অনুক্রমের দৈর্ঘ্যের উপর সাধারণ আরোহ ব্যবহার
করে আমরা সূত্র সম্বন্ধে উপপাদ্য প্রমাণ করতে পারি।

\begin{lem}
\ollabel{lem:fseq-init}
ধরি $\tuple{!A_0,\dotsc,!A_n}$ হলো~$!A_n$-এর একটি গঠন-অনুক্রম এবং
$k \leq n$। তাহলে $\tuple{!A_0,\dotsc,!A_k}$ হলো~$!A_k$-এর একটি
গঠন-অনুক্রম।
\end{lem}

\begin{proof}
অনুশীলনী।
\end{proof}

\begin{thm}
\ollabel{thm:fseq-frm-equiv}
$\Frm[L_0]$ হলো $\Lang L_0$ ভাষায় এমন সব প্রতীকক্রমের সেট,
যাদের গঠন-অনুক্রম আছে।
\end{thm}

\begin{proof}
$\Lang L_0$ ভাষায় গঠন-অনুক্রম আছে এমন সব প্রতীকক্রমের সেটকে $F$ ধরি।
\olref[pl][syn][fseq]{prop:formed}-এ আমরা দেখেছি যে
$\Frm[L_0] \subseteq F$; এবার বিপরীত অন্তর্ভুক্তি প্রমাণ করি।

ধরি $!A$-এর গঠন-অনুক্রম $\tuple{!A_0,\dotsc,!A_n}$।
~$n$-এর উপর প্রবল আরোহ ব্যবহার করে আমরা প্রমাণ করি যে
$!A \in \Frm[L_0]$। আমাদের আরোহ-অনুমান হলো: শেষ সূচক $m < n$ এমন
গঠন-অনুক্রমবিশিষ্ট প্রতিটি প্রতীকক্রম~$\Frm[L_0]$-এর অন্তর্গত।
গঠন-অনুক্রমের সংজ্ঞা অনুসারে, হয় $!A_n$ পরমাণু, নয়তো এমন $j,k < n$
অবশ্যই আছে যাতে নিচের কোনো একটি সত্য:
\begin{enumerate}
    \tagitem{prvNot}{$!A_n \ident \lnot !A_j$।}{}%
    \tagitem{prvAnd}{$!A_n \ident (!A_j \land !A_k)$।}{}%
    \tagitem{prvOr}{$!A_n \ident (!A_j \lor !A_k)$।}{}%
    \tagitem{prvIf}{$!A_n \ident (!A_j \lif !A_k)$।}{}%
    \tagitem{prvIff}{$!A_n \ident (!A_j \liff !A_k)$।}{}%
\end{enumerate}
এখন ক্ষেত্রভেদে যুক্তি দিই। $!A_n$ পরমাণু হলে
$!A_n \in \Frm[L_0]$। অন্য দিকে, ধরি $!A_n \ident
(!A_j \land !A_k)$। \olref[pl][syn][fseq]{lem:fseq-init} অনুসারে
$\tuple{!A_0,\dotsc,!A_j}$ ও $\tuple{!A_0,\dotsc,!A_k}$ যথাক্রমে
$!A_j$ ও~$!A_k$-এর গঠন-অনুক্রম। এগুলি~$!A$-এর গঠন-অনুক্রমের প্রকৃত
প্রারম্ভিক উপ-অনুক্রম বলে উভয়ের শেষ সূচক~$n$-এর চেয়ে ছোট। তাই
আরোহ-অনুমান থেকে $!A_j$ ও~$!A_k$ উভয়ই~$\Frm[L_0]$-এর অন্তর্গত;
সুতরাং !!a{formula}-এর সংজ্ঞা অনুসারে
$(!A_j \land !A_k)$-ও তাই। অন্য ক্ষেত্রগুলিও একই ধরনের যুক্তি থেকে আসে।
% BN-SRC-020 TARGET-CORRECTION-BEGIN
\par\smallskip\noindent\textnormal{\small\textbf{উৎস-সংশোধন (BN-SRC-020):} হিমায়িত ইংরেজি উৎসে এই ক্ষেত্রটির শুরুতে $!A \equiv (!A_j \land !A_k)$ আছে। কিন্তু সংজ্ঞা ও ঠিক আগের ক্ষেত্রতালিকা গঠন-অনুক্রমের শেষ পদ~$!A_n$ এবং সংকেতবিন্যাসগত অভিন্নতার সংকেত~$\ident$ ব্যবহার করে। বাংলা পাঠে তাই $!A_n \ident (!A_j \land !A_k)$ করা হয়েছে।} % BN-SRC-020 TARGET-CORRECTION-END
\end{proof}

\end{document}
